\begin{recipe}{Mini-aardappelquiches}
\kicker[Voorgerecht,Vlees,Frans,Kerst]{Voorgerecht · vlees · Frans · Kerst}
\index[register]{Mini-aardappelquiches}
\index[register]{Aardappelen}
\index[register]{Spek}
\index[register]{Parmigiano Reggiano}

\meta{40 minuten \dvd 12 stuks \dvd Oven 200\,°C}

\lettrine{M}{ini-aardappelquiches} met een knapperige aardappelbodem en een vulling
van spek, bosui en Parmezaan. Een paar keer gemaakt voor het kerstdiner op de
basisschool. Van tevoren te maken en ook koud lekker.

\blockrule{Ingrediënten}
\begin{ingredients}
  \ing{600 g}{kruimige aardappelen, geschild}\index[register]{Aardappelen}
  \ing{100 g}{ontbijtspek, in stukjes}\index[register]{Spek}
  \ing{3 el}{zonnebloemolie}
  \ing{3}{bosuien, in ringetjes}
  \ing{75 g}{Parmigiano Reggiano, geraspt}\index[register]{Parmigiano Reggiano}
  \ing{125 ml}{verse slagroom}
  \ing{3}{eieren}
  \ing{2 tl}{gedroogde rozemarijn}
  \ingb{zout}
  \ingb{snufje nootmuskaat of cayennepeper}
\end{ingredients}

\blockrule{Bereiding}
\tip{Knijp de geraspte aardappel goed uit voor je hem in de vorm doet. Minder vocht geeft een knapperigere bodem. Het spek is te vervangen door vegetarisch rookvlees of weg te laten.}
\begin{steps}
  \item Verwarm de oven voor op 200\,°C (boven- en onderwarmte). Vet een muffinvorm met 12 holtes in.
  \item Rasp de aardappelen fijn en meng ze meteen met de olie: zo sluit het vet de aardappel af en verkleurt hij niet. Meng er de rozemarijn en een snuf zout door en knijp het vocht eruit.
  \item Verdeel het aardappelmengsel over de holtes en druk goed aan tot er bakjes
        ontstaan. Bak ca.\ 15 minuten tot goudbruin. Haal uit de oven.
  \item Bak het spek op middelhoog vuur in een droge koekenpan in 4--6 minuten krokant:
        het spek laat zelf genoeg vet los. Klop de slagroom,
        de geraspte kaas, de eieren en een snufje nootmuskaat of cayennepeper samen
        met een garde.
  \item Verdeel het spek en het grootste deel van de bosuien over de aardappelbakjes,
        en houd een handjevol bosui achter. Schenk het ei-roommengsel erover.
  \item Verlaag de oventemperatuur naar 175\,°C en bak de quiches nog ca.\ 15 minuten
        gaar. Laat iets afkoelen en lepel ze met een lepeltje uit de vorm. Strooi de
        achtergehouden bosuiringetjes erover, dan heb je kleur op het goudbruine bakje.
\end{steps}
\end{recipe}
